\begin{abstract}
A reporting rule should be evaluated against the quantity and uncertainty it
claims to report. We audit a 50-cell rule in two colorectal single-cell cohorts
that had been justified by point-estimate recovery. In 5{,}000 simulations per
setting, its nominal 95\% percentile intervals reject a true null in
9.8--32.2\% of studies at the reporting boundary. A paired comparison with
Welch intervals improves effect coverage but leaves null rejection at
8.9--30.8\%; changing the interval alone does not repair the rule. We also
show how a sampled reference inflates apparent coverage, explain why sparse
samples can produce misleadingly high detection, and document grid-dependent
threshold selection. A synthetic external-control example illustrates a
related mistake: a reference for the wrong population favors an estimator
of the wrong effect. The practical lesson is to specify the target, test
null and alternative performance together, and evaluate an interval and its
reporting threshold as a combined rule.
\end{abstract}
